\chapter*{Kurzfassung}\addcontentsline{toc}{chapter}{Kurzfassung}
Die Hochabtastung von Punktwolken erhöht die Anzahl räumlicher Abtastpunkte.
Ein Nutzen für nachgelagerte Wahrnehmungsaufgaben lässt sich jedoch nicht
allein aus geometrischen Abständen ableiten. Diese Arbeit untersucht die
Verdichtung einer ursprünglichen Eingabe sowie die Rekonstruktion nach
einer Reduktion auf ein Viertel der Punkte. Auf ModelNet40 werden EAR,
PDANS, PU-Net, PU-GCN und PU-EdgeFormer mit separat trainierten
PointNet++-Klassifikatoren und geometrischen Referenzen gleicher
Punktanzahl verglichen. Auf KITTI werden PointRCNN und CenterPoint
untersucht. Eine PU-GCN-Eingabe- und Detektoranpassungsstudie über drei
Epochen umfasst 20 Auswertungen. Anschließend wird das Training mit
Eingaben unter Erhalt der beobachteten Punkte verlängert. Jede
Hauptauswertung verwendet 3.769 Validierungsframes.

Bei der protokollierten, über Batches gemittelten Genauigkeit übertrifft
keine getestete Methode in ModelNet40-Linie~A die beste
Gesamtgenauigkeit der ursprünglichen Darstellung mit 1.024 Punkten
von 91,95\%. Nach der Ausdünnung verbessert PU-Net diese Genauigkeit
von 90,85\% auf 91,27\%, während PU-GCN in beiden Linien den niedrigsten
mittleren Chamfer-Abstand erreicht. Auf KITTI verschlechtern direkte
PU-GCN-Eingaben die Ergebnisse der unveränderten Detektoren. Der Erhalt
gemessener Punkte und die Anpassung der Detektoren verbessern die
Ergebnisse, übertreffen jedoch nicht die jeweiligen Referenzen. Das
verlängerte Training verbessert alle vier Detektor-Linien-Kombinationen
weiter und erfüllt das festgelegte Plateaukriterium für die
Car-Validierungsmetrik. Dennoch verbleiben Rückstände von 3,31 bis
11,28 AP-Punkten gegenüber den jeweiligen Linienreferenzen. Das
PU-GCN-Netz selbst verwendet weiterhin seine veröffentlichten Gewichte.

Diagnostische Versuche zu Patch-Normalisierung, Punktverarbeitung und
Voxelbelegung zeigen, warum gespeicherte Punktdichte und wirksame
Modelleingabe verschieden sind. Die Aussagekraft wird durch nur einen
Trainingsseed, die Auswahl des besten Klassifikators anhand des Testsatzes,
ungleiche Trainingsbudgets der Detektoren und die Auswahl ihrer
Checkpoints anhand desselben Validierungssatzes begrenzt. Geometrische Qualität allein reicht somit nicht aus, um einen Nutzen
für nachgelagerte Aufgaben nachzuweisen.
